\documentclass[10pt,a4paper]{amsart}
\usepackage[margin=19mm]{geometry}
\usepackage{amsmath,amssymb,mathtools}
\usepackage[scr=rsfs,cal=euler]{mathalfa}
\usepackage{xfrac}
\usepackage[unicode,hidelinks,bookmarks=false]{hyperref}
\newcommand{\ie}{i.e.\ }
\newcommand{\cf}{cf.\ }
\newcommand{\resp}{respectively\ }
\newcommand{\Subsection}{\subsection}
\providecommand{\nobreakdash}{\nobreak\hbox{-}}
\setlength{\emergencystretch}{2em}
\multlinegap0pt
\usepackage{bm}
\usepackage{bbm}
\usepackage{dsfont}
\usepackage{xspace}
\usepackage{cancel}
\usepackage{mathtools}
\usepackage{amsthm}
\makeatletter
\@ifundefined{theorem}{\newtheorem{theorem}{Theorem}}{}
\@ifundefined{proposition}{\newtheorem{proposition}[theorem]{Proposition}}{}
\makeatother
\DeclareMathOperator{\lat}{lat}
\newcommand{\bigfree}{\mathop{\ast}}
\newcommand{\one}{\mathbf{1}}
\title[Free Products of Banach Lattices]{Free Products of Banach Lattices}
\author{Gonzalo Mart\'inez-Fern\'andez, Pedro Tradacete}
\date{Source: arXiv:2605.28988v1 (CC BY 4.0)}
\begin{document}
\maketitle

\noindent\emph{Source and licence.} Free Products of Banach Lattices. Version arXiv:2605.28988v1 (CC BY 4.0), distributed under the Creative Commons Attribution 4.0 International licence, as recorded in the arXiv licence metadata for this exact version and shown on the abstract page https://arxiv.org/abs/2605.28988v1. Full rights notice: licenses/arxiv-2605.28988-CC-BY-4.0.txt.\par\smallskip

\noindent\textit{Editorial scope.} Counted unit: the Theorem whose source label is \texttt{thm:CK-join} in the article (source0.tex lines 937--950, source label \texttt{thm:CK-join}), delivered together with the article's own complete original proof of it (source0.tex lines 952--1073, one \texttt{proof} environment of 122 lines). No mathematics is written, repaired or re-derived: statement and proof are the article's own text line for line. Required context retained uncounted: prop:maximal-norm (lines 320--336); prop:strong-unit (lines 887--893). Every definition, display and label of those lines is retained unchanged. Omitted from this package and not redistributed: 1--319, 337--886, 894--936, 951--951, 1074--2004. Nothing inside a retained excerpt is omitted or altered: every retained body is delivered exactly as the article's lines are written, so no omission receipt of any kind accompanies this package. Citations: the retained excerpts carry no citation command at all, so no bibliography entry of the article is transcribed and no reference is redistributed. Reference adaptation: none was needed and none was made; every reference inside the delivered unit resolves inside it, so no unresolved reference or citation is produced. Printed slips kept literal: no printed word, symbol, hypothesis or number of the counted statement or of its proof was altered. Layout and class: the article's own class is \texttt{amsart} with packages including amsmath, amssymb, amsthm, url, a4wide, graphicx, float, mathrsfs, bbold, xcolor, xy; the wrapper is the lane's pinned \texttt{amsart} editorial wrapper at 10pt on A4, which restates the theorem-like environments the retained text needs and adds the article's own macro definitions that the retained text needs (\texttt{\textbackslash bigfree}, \texttt{\textbackslash lat}, \texttt{\textbackslash one}), with extra wrapper packages (bm, bbm, dsfont, xspace, cancel, mathtools). The delivered numbering is the unit's own. Assets: no retained span contains a figure environment, a graphics-inclusion command, a PDF-page inclusion, a table or a TikZ picture; no image file is copied into this package, the package declares no asset dependency and no image file is redistributed with it. Licence: version arXiv:2605.28988v1 of this article is distributed under the Creative Commons Attribution 4.0 International licence, as recorded in the arXiv licence metadata for this exact version and shown on the abstract page https://arxiv.org/abs/2605.28988v1. A full rights notice is delivered at licenses/arxiv-2605.28988-CC-BY-4.0.txt. Characters: every retained excerpt is pure ASCII, so the delivered engine, XeLaTeX with its default Latin Modern text font, reports no missing character.
\begin{proposition}\label{prop:maximal-norm}
Let
\[
Z=\lat\{\varphi_i(X_i):i\in I\}\subseteq \bigfree_{i\in I}X_i.
\]
Suppose that $\|\cdot\|$ is a lattice norm on $Z$ such that
\[
\|\varphi_i(x)\|\leq \|x\|_{X_i}
\qquad (x\in X_i,\ i\in I).
\]
Then
\[
\|z\|\leq \|z\|_{\ast}
\qquad (z\in Z),
\]
where $\|\cdot\|_{\ast}$ denotes the free product norm.
\end{proposition}

\begin{proposition}\label{prop:strong-unit}
The element
\[
e=\varphi_1(\one_{K_1})+\varphi_2(\one_{K_2})\in C(K_1)\ast C(K_2)
\]
is a strong order unit.
\end{proposition}

\begin{theorem}\label{thm:CK-join}
Let $K_1$ and $K_2$ be compact Hausdorff spaces. The spaces $C(K_1)\ast C(K_2)$ and $C(K_1\ast K_2)$ are lattice isomorphic (with constant 2). Under this identification, the canonical copies of $C(K_1)$ and $C(K_2)$ are given by
\[
\varphi_1(g)=L(g,0),
\qquad
\varphi_2(h)=L(0,h).
\]
and the free product norm of a function $f\in C(K_1*K_2)$ is given by
\[
\|f\|_*
=\inf\Big\{a+b:\ a,b\in\mathbb R_+,\ |f[x,y,t]|\leq (1-t)a+tb\ \text{for all } [x,y,t]\in K_1\ast K_2\Big\}.
\]

\end{theorem}

\begin{proof}
By Proposition~\ref{prop:strong-unit}, the free product $C(K_1)\ast C(K_2)$ has a strong order unit. By Kakutani's representation theorem, there exists a compact Hausdorff space $F$ such that
\[
C(K_1)\ast C(K_2)\simeq C(F),
\]
where \[
F=\{u:C(K_1)*C(K_2)\to\mathbb R:u\text{ is a lattice homomorphism and } u(e)=1\}
\]
with the product topology and $e=\varphi_1(\bold 1_{K_1})+\varphi_2(\bold 1_{K_2})$ is the strong unit from Proposition \ref{prop:strong-unit}. We claim that $F\cong K_1*K_2$.

Let
\[
\pi:K_1\times K_2\times [0,1]\to K_1\ast K_2
\]
be the quotient map. 

Let $u\in F$, and set $u_i=u\circ \varphi_i$ for $i=1,2$. Each $u_i$ is a lattice homomorphism from $C(K_i)$ to $\mathbb{R}$, so $u_i=c_i\delta_{x_i}$ for some $c_i\geq 0$ and $x_i\in K_i$. If $c_i=0$, we can choose $x_i\in K_i$ arbitrarily. Since
\[
1=u(e)=u_1(\one_{K_1})+u_2(\one_{K_2})=c_1+c_2,
\]
we may define
\[
\Phi(u)=[x_1,x_2,c_2]\in K_1\ast K_2.
\]
The identifications of points in the join ensure that $\Phi$ is well defined when $c_2=0$ or $c_2=1$.

The map $\Phi$ is bijective. Indeed, given $[x_1,x_2,t]\in K_1\ast K_2$, the lattice homomorphisms $(1-t)\delta_{x_1}$ on $C(K_1)$ and $t\delta_{x_2}$ on $C(K_2)$ extend uniquely to a lattice homomorphism
\[
u:C(K_1)\ast C(K_2)\to \mathbb{R}
\]
with $u(e)=1$, and $\Phi(u)=[x_1,x_2,t]$. Conversely, if $\Phi(u)=\Phi(v)=[x_1,x_2,t]$, then
\[
u_1=(1-t)\delta_{x_1}=v_1,
\qquad
u_2=t\delta_{x_2}=v_2,
\]
and uniqueness in the free product gives $u=v$.

It remains to prove continuity. Let $(u^\alpha)$ be a net in $F$ converging to $u\in F$. Write
\[
u_i^\alpha=c_i^\alpha\delta_{x_i^\alpha},
\qquad
u_i=c_i\delta_{x_i}
\]
whenever the corresponding functional is non-zero, and choose $x_i^\alpha\in K_i$ arbitrarily if $u_i^\alpha=0$. Then
\[
c_i^\alpha=u_i^\alpha(\one_{K_i})\longrightarrow u_i(\one_{K_i})=c_i
\qquad (i=1,2).
\]
If $0<c_i<1$, then eventually $c_i^\alpha>0$, and for every $f\in C(K_i)$ we have
\[
c_i^\alpha f(x_i^\alpha)=u_i^\alpha(f)\longrightarrow u_i(f)=c_i f(x_i).
\]
Since $c_i^\alpha\to c_i>0$, it follows that $f(x_i^\alpha)\to f(x_i)$ for every $f\in C(K_i)$, hence $x_i^\alpha\to x_i$. Therefore
\[
\Phi(u^\alpha)\to \Phi(u)
\]
whenever $0<c_2<1$.

If $c_2=0$, then $c_1=1$ and $x_1^\alpha\to x_1$ by the same argument applied to $u_1^\alpha$. A basic neighborhood of $\Phi(u)$ in $K_1\ast K_2$ has the form
\[
\pi(U\times K_2\times [0,\varepsilon))
\]
for some neighborhood $U$ of $x_1$ in $K_1$. Since $x_1^\alpha\to x_1$ and $c_2^\alpha\to 0$, we eventually have
\[
\Phi(u^\alpha)\in \pi(U\times K_2\times [0,\varepsilon)).
\]
The case $c_2=1$ is analogous. Thus $\Phi$ is continuous.

Since $F$ is compact and $K_1\ast K_2$ is Hausdorff, $\Phi$ is a homeomorphism. Therefore
\[
C(K_1)\ast C(K_2)\simeq C(F)\simeq C(K_1\ast K_2).
\]

To see what the canonical inclusions look like under this identification, fix $f\in C(K_1)$ and $[x,y,t]\in K_1\ast K_2$. Let
\[
u_{x,y,t}:C(K_1)\ast C(K_2)\to \mathbb{R}
\]
be the lattice homomorphism corresponding, via the argument above, to the point $[x,y,t]$. Then
\[
\varphi_1(f)[x,y,t]
=u_{x,y,t}(\varphi_1(f))
=((1-t)\delta_x\ast t\delta_y)(\varphi_1(f))
=(1-t)f(x)
=L(f,0)[x,y,t].
\]
The formula for $\varphi_2(g)$ is identical.

Lastly, we compute the free product norm. Define for every $f\in C(K_1*K_2)$ the following lattice seminorm
\[
p(f)=\inf\Big\{a+b:\ a,b\in\mathbb R_+,\ |f|\leq L(a,b)\Big\}.
\]
This is a lattice norm, and it is equivalent to the supremum norm because
\[
\|f\|_\infty\leq p(f)\leq 2\|f\|_\infty.
\]

If $|f|\leq L(a,b)$, then
\[
L(a,b)=L(a,0)+L(0,b)=a\,\varphi_1(\one_{K_1})+b\,\varphi_2(\one_{K_2}),
\]
and therefore
\[
\|f\|_{\ast}\leq a\|\varphi_1(\one_{K_1})\|_{\ast}+b\|\varphi_2(\one_{K_2})\|_{\ast}=a+b.
\]
Taking the infimum yields $\|f\|_{\ast}\leq p(f)$.

On the other hand, if $g\in C(K_1)$ then $|L(g,0)|\leq L(a,b)$ implies $|g(x)|\leq a$ for every $x\in K_1$, hence
\[
p(L(g,0))\leq\|g\|_\infty.
\]
Similarly,
\[
p(L(0,h))\leq\|h\|_\infty
\qquad (h\in C(K_2)).
\]
Thus $p$ agrees with the original norms on the two canonical copies. By Proposition~\ref{prop:maximal-norm},
\[
p(f)\leq \|f\|_{\ast}
\]
on the dense sublattice generated by those copies, and therefore on all of $C(K_1\ast K_2)$ by continuity. Hence $p=\|\cdot\|_{\ast}$.
\end{proof}

\end{document}
